% partes/parte1/intro.tex
\part{Comportamiento, perfil y gestión de datos}

\chapter*{Introducción a la Parte I}
\addcontentsline{toc}{chapter}{Introducción a la Parte I}

Un mismo clic significa cosas distintas según quién lo interprete. Para el
consumidor, es el cierre de una decisión: comparó, dudó, se decidió. Para la
plataforma que registró ese clic, es un dato entre millones, destinado a
alimentar un modelo que predice el próximo clic de alguien más. Esta Parte
trata sobre esa doble naturaleza del consumo digital: es, a la vez, una
experiencia vivida por una persona y un rastro estructurado que una
organización puede analizar, agregar y usar para decidir.

El programa oficial de la unidad llama a esta Parte \enquote{Comportamiento, perfil y
gestión de datos del consumidor digital}, y su competencia es identificar la
evolución del consumidor digital a partir de sus tipos de perfil y de la
gestión de los datos que genera. Cuatro capítulos la desarrollan. El primero
sitúa históricamente al consumidor digital: qué cambió, qué no, y por qué
conviene desconfiar de las periodizaciones ordenadas en versiones (\enquote{del
consumidor 1.0 al 5.0}) que la literatura de consultoría presenta como si
fueran hechos científicos. El segundo aborda el perfil y la segmentación:
cómo se agrupa a los consumidores, y por qué la pertenencia rígida a un solo
segmento es, en muchos mercados, una simplificación conveniente más que una
descripción fiel del comportamiento real. El tercero entra en la gestión de
datos propiamente dicha: cómo se crea valor a partir de ellos, qué significa
\enquote{escuchar} a un consumidor digital, y qué obligaciones legales y éticas trae
consigo recolectar esa información en México en 2026. El cuarto cierra con la
base tecnológica que hace posible todo lo anterior: buscadores, plataformas,
redes sociales y comercio electrónico, tratados no como un catálogo de marcas
sino como un conjunto de criterios de análisis que sobrevivan al cambio de
esas marcas.

Un hilo conecta los cuatro capítulos: la gestión de datos del consumidor
digital no es neutral. Cada decisión de qué medir, cómo segmentar y qué
inferir tiene consecuencias sobre la privacidad, la autonomía y, en última
instancia, la confianza del consumidor en la organización que lo observa. Por
eso cada capítulo cierra con un recuadro de ética y responsabilidad con los
datos: no como formalidad, sino porque la pregunta \enquote{¿debería medirse esto?}
es tan parte del oficio del mercadólogo digital como la pregunta \enquote{¿cómo se
mide esto?}.

Las prácticas oficiales de esta Parte —\enquote{El consumidor digital y su perfil} y
\enquote{Buscadores y plataformas en la gestión de datos del consumidor digital}— se
integran en las secciones de laboratorio de cada capítulo y se consolidan al
final de la Parte, en la sección \enquote{Prácticas oficiales de la unidad}.
